% !PRINTED 186
\weberpage{186}
Denn setzen wir

\[
D = \begin{vmatrix} a_1 - 1, & a_2, & a_3\\ b_1, & b_2 - 1, & b_3\\ c_1, & c_2, & c_3 - 1 \end{vmatrix},
\]

so erhalten wir aus den Relationen (4):

\[
|A|\, D = - D,
\]

also, wenn (5) besteht, $D = 0$. Demnach lassen sich drei
Gr\"ossen $\lambda, \mu, \nu$ aus den Gleichungen:

\weq{6}{\begin{array}{l}
\lambda = a_1 \lambda + a_2 \mu + a_3 \nu\\
\mu = b_1 \lambda + b_2 \mu + b_3 \nu\\
\nu = c_1 \lambda + c_2 \mu + c_3 \nu\\
\lambda^2 + \mu^2 + \nu^2 = 1
\end{array}}

bestimmen, und diese drei Gr\"ossen $\lambda, \mu, \nu$ bestimmen die Richtung einer geraden Linie, die mit den Axen $y_1, y_2, y_3$ dieselben
Winkel einschliesst, wie mit den Axen $x_1, x_2, x_3$. Eine in dieser
Richtung durch den Coordinatenanfangspunkt gelegte Linie ist
die Drehungsaxe.

In dem Falle $|A| = - 1$ trifft dieser Schluss nicht mehr zu.

Daneben besteht noch eine andere Deutung der orthogonalen
Substitution $(y) = A (x)$, wonach $(x)$ und $(y)$ die Coordinaten
zweier verschiedener Punkte $x$ und $y$ sind, bezogen auf ein und
dasselbe Coordinatensystem. Wenn $x$ einen Raumtheil (Linie,
Fl\"ache oder K\"orper) \"uberstreicht, so erf\"ullt der entsprechende
Punkt $y$ einen congruenten Raumtheil (wenn $|A| = - 1$ ist, einen
spiegelbildlich gleichen). Dies wird durch die folgende Betrachtung dargethan.

Die Gruppe der \gesperrt{eigentlichen} orthogonalen Substitutionen
ist \"aquivalent mit einer Gruppe, die man aus den verschiedenen
Stellungen eines um einen festen Punkt drehbaren K\"orpers bilden
kann. Diese Gruppe erh\"alt man, wenn man eine beliebige Stellung $E$ als Einheit annimmt, aus der man in irgend eine andere
Stellung $A$ gelangt durch Drehung um eine bestimmte Axe mit
einem bestimmten Winkel. Ist $B$ eine dritte Stellung, so hat
man unter der zusammengesetzten Stellung $AB$ die Stellung zu
verstehen, die man erh\"alt, wenn man die Drehung, die zu $A$
gef\"uhrt hat, nicht von der Einheitsstellung, sondern von der
Stellung $B$ aus vollf\"uhrt. Denn nimmt man ein mit dem K\"orper
in starrer Verbindung stehendes Axensystem an, z. B. das System
der Haupttr\"agheitsaxen, und bezeichnet mit $(x)$ die Coordinaten
%CONT>>
